\fsection{Viabilidad de los proyectos propuestos}
Se presentan dos propuestas de proyecto viables técnica y económicamente. La propuesta A es la que se desarrolla en este proyecto; la propuesta B se plantea como ampliación sobre ella:

\fsubsection{
	Propuesta A — Sistema de bin picking en entorno no estructurado:
}
El robot extrae piezas de un contenedor de entrada desordenado
mediante visión artificial e IA y las deposita en la zona de salida, sin necesidad de posicionarlas previamente. Aplicable en
líneas de producción y centros de distribución donde las piezas llegan a granel. Viabilidad técnica alta: el sistema base de
Siemens está validado y se dispone de todo el equipamiento en el centro, por lo que no requiere inversión adicional.

\fsubsection{
	Propuesta B — Célula de e-Commerce simulada con identificación de artículos:
}

El robot gestiona el picking de pedidos desde un contenedor de entrada,
identifica cada artículo (por RFID o código QR) y lo deposita en la bandeja de salida correspondiente al pedido, simulando el flujo
de un almacén de comercio electrónico. Esta propuesta añade mayor complejidad en la lógica de clasificación y en la gestión
de múltiples destinos, y es directamente extrapolable a un entorno real de operador logístico. Requiere adquirir el
sistema de identificación, por lo que se deja como trabajo futuro.
